% 01-uvod.tex — Introduction + Related Work (TASK §F.2).
\section{Uvod}

Klasično programiranje robota zadaje putanju unaprijed. Fizička interakcija
čovjeka i robota (\emph{physical Human--Robot Interaction}, pHRI) mijenja tu
podjelu jer čovjek vlastitim pokretom neposredno određuje što robot radi, u
trenutku u kojem se pokret događa. Time nestaje razdvojenost operatera i stroja
koja je dosad bila glavni sigurnosni mehanizam, pa se sigurnost mora preseliti u
sam upravljački lanac. Problem je i mjerni. Ljudski pokret nosi šum, tremor i
prekide u vidljivosti, a robot traži naredbe koje su glatke, ograničene po
brzini i predvidljive u svakom taktu.

Zadatak koji stoji pred ovim radom postavljen je konkretno. Treba izgraditi
sustav koji u stvarnom vremenu prima položaje markera ljudske ruke iz sustava
OptiTrack, iz njih računa zakret odabranog ljudskog zgloba, taj zakret
preslikava na jedan zglob robota Universal Robots, naredbu šalje robotu mrežom i
radi kontinuirano, uz jasno određene granice sigurnog ponašanja. Odabrani par
zglobova anatomski je podudaran, pri čemu kut fleksije ljudskog lakta vodi lakat robota,
zglob J3. Savijanje ruke savija robota, pa je veza između pokreta i posljedice
očita i bez objašnjenja. Sustav opisan ovdje ispunjava taj zahtjev i proširuje
ga. Istim obradnim lancem i u istom taktu sustav vodi svih šest zglobova kolaborativnog
robota UR3e iz pokreta jedne ljudske ruke, pri čemu je obvezni kanal lakta na J3
jedan od tih šest zglobova, izmjeren istim alatom i u svakom pokusu kao i preostalih pet.

Optički sustavi za praćenje pokreta, poput sustava OptiTrack, mjere položaj
segmenata tijela s milimetarskom točnošću na frekvenciji od 120 Hz. Premda je to
zadovoljavajuće za teleoperaciju, kolaborativni robot UR3e ima tvorničku ponovljivost
pozicioniranja od $\pm 0{,}03$\,mm, pa optički sustav i šum markera predstavljaju
glavni limitator točnosti cjelokupnog postava. Izazov koji nije riješen samim
mjerenjem jest put od optičke rekonstrukcije do pogona robota, odnosno pitanje koliko se kašnjenja
nakupi duž obradnog lanca, koliki udio dolazi od filtriranja, a koliki od dinamike
samog robota, te pri kojem iznosu kašnjenja operater gubi osjećaj izravnog upravljanja.
Ovaj rad eksperimentalno istražuje taj put i mjeri ga na stvarnom kolaborativnom
robotu u laboratoriju.

\subsection{Arhitektura sustava}

Sustav je izveden kao obradni lanac neovisan o izvoru podataka i odredištu
naredbi. Isti kod radi s reproduciranom snimkom i sa živim hardverom, a zamjena
izvora (snimka ili živi tok podataka protokolom NatNet) i odredišta
(simulator URSim ili stvarni UR3e) provodi se promjenom konfiguracije. Tok podataka
prikazuje slika~\ref{fig:arhitektura}, a proteže se od izvora i prijemnika preko
izračuna zglobnih kanala, kauzalnog filtra i sigurnosnog sloja do robota
putem sučelja RTDE. Uz svaki takt bilježe se vremenske oznake te, odvojeno,
svi kutovi (ulazni, filtrirani, ciljni, naredbeni i stvarno izmjereni kut robota).
Ta zadnja stavka razlikuje mjerenje vjernosti od mjerenja brzine, i o njoj ovisi
najveći dio rezultata u poglavlju~\ref{sec:rezultati}.

\begin{figure*}[t]
  \centering
  \begin{tikzpicture}[
      font=\footnotesize,
      box/.style={draw, rounded corners=2pt, align=center, minimum height=8mm,
                  inner xsep=4pt, text width=20mm},
      arr/.style={-{Latex[length=2mm]}, thick},
    ]
    \node[box] (src)    {Izvor\\\scriptsize OptiTrack ili snimka (UDP, NatNet)};
    \node[box, right=5mm of src]  (rx)   {Prijem\\\scriptsize gubitak, redoslijed, frekvencija};
    \node[box, right=5mm of rx]   (ang)  {Zglobni kanali\\\scriptsize poze krutih tijela};
    \node[box, right=5mm of ang]  (filt) {Kauzalni filtar\\\scriptsize One-Euro po zglobu};
    \node[box, right=5mm of filt] (safe) {Sigurnosni sloj\\\scriptsize brzina, raspon, zadržavanje};
    \node[box, right=5mm of safe] (rob)  {Robot\\\scriptsize RTDE (UR3e / URSim)};
    \draw[arr] (src) -- (rx);
    \draw[arr] (rx) -- (ang);
    \draw[arr] (ang) -- (filt);
    \draw[arr] (filt) -- (safe);
    \draw[arr] (safe) -- (rob);
    \draw[decorate, decoration={brace, amplitude=4pt, mirror}]
      ([yshift=-3mm]src.south west) -- ([yshift=-3mm]rob.south east)
      node[midway, below=2pt] {\scriptsize po taktu se bilježe vremenske oznake i svi kutovi, uključujući stvarni kut robota};
  \end{tikzpicture}
  \caption{Obradni lanac neovisan o izvoru i odredištu. Zamjena izvora (snimka
  ili živi tok podataka) i odredišta (URSim ili UR3e) je promjena
  konfiguracije.}
  \label{fig:arhitektura}
\end{figure*}

Rad daje tri doprinosa. Prvi je cjelovit upravljački lanac s malim računskim kašnjenjem,
od optičkog mjerenja do stvarnog kolaborativnog robota, sa sigurnosnim slojem
čije djelovanje ostaje zapisano u teleometriji. Drugi doprinos je mjerni, kroz
analitičko i eksperimentalno razlaganje ukupnog kašnjenja ruka--robot na filtar, sigurnosni sloj i regulaciju pogona,
čime se egzaktno pokazuje koji dio lanca određuje dinamički odziv. Treći je vrednovanje
sustava u području fizičke interakcije čovjeka i robota, gdje se objektivne mjere
odziva povezuju s opažanjima operatera pri umjetno unesenom kašnjenju i analizi
granica upravljivosti.

\subsection{Pregled povezanih radova}

\textbf{Sigurnost u pHRI-u.} Norma ISO/TS 15066 definira način rada s
ograničenjem snage i sile (\emph{power-and-force limiting}, PFL), u kojem se
najveća dopuštena brzina robota izvodi iz efektivne mase, dopuštene sile i
krutosti zahvaćenog dijela tijela. Ghanbarzadeh i Najafi na toj osnovi predlažu
regulator s promjenjivom impedancijom za siguran kontakt
\cite[str.~4]{ghanbarzadeh2025}. Peng i suradnici razvijaju dinamičko praćenje
brzine i razdvojenosti (\emph{speed and separation monitoring}, SSM) u kojem se
zaštitna udaljenost, a time i brzina robota, prilagođavaju blizini čovjeka
\cite[str.~1]{peng2025}. Oba pristupa motiviraju naš sigurnosni sloj. Budući da
robotom upravlja čovjek koji stoji uz njega, umjesto pune impedancije
primjenjujemo jednostavnije, ali deterministički provjerljivo ograničenje brzine
i raspona, čije granice odjeljak~\ref{sec:postav} provjerava PFL proračunom;
poglavlje~\ref{sec:rezultati} pokazuje koliko je često taj sloj doista djelovao.

\textbf{Teleoperacija, kašnjenje i osjećaj kontrole.} Wang i suradnici pokazuju
da pod vremenski promjenjivim mrežnim kašnjenjem postoji kompromis između
stabilnosti i transparentnosti bilateralne teleoperacije
\cite[str.~5]{wang2025}. Louca i suradnici eksperimentalno potvrđuju da porast
latencije sustavno pogoršava uspješnost, točnost i kontaktne sile pri
teleoperaciji \cite[str.~14]{louca2024}. Shi i suradnici optimiraju latenciju i
pouzdanost mrežne teleoperacije s povratom sile \cite[str.~9]{shi2025}, dok
Zainudin i suradnici smanjuju kašnjenje upravljanja robotom UR10e
\textit{multi-thread} obradom, izbjegavajući inverznu kinematiku i
opsežnu kalibraciju \cite[str.~2]{zainudin2025}. Ti radovi određuju što mjerimo
i kako to tumačimo. Naš je put lokalan, pa je red veličine brži od njihovih
mrežnih scenarija, zbog čega mjereno kašnjenje ovdje ne dolazi
od mreže, nego od filtriranja signala.

\textbf{Preslikavanje pokreta na robota.} Weigend i suradnici procjenjuju pozu
ruke iz oskudnih senzora za sveprisutno upravljanje robotom uz mjeru
nesigurnosti \cite[str.~1]{weigend2023}, a okvir iRoCo probabilističkim
diferencijabilnim filtrima uravnotežuje preciznost upravljanja i slobodu
kretanja korisnika \cite[str.~1]{weigend2024}. Za razliku od tih pristupa
zasnovanih na učenju i procjeni cijele poze, naše je preslikavanje izravno i
eksplicitno jer šest kanala izvedenih iz poza triju krutih tijela vodi šest
zglobova. Time gubimo općenitost, a dobivamo lanac u kojem je svaka
transformacija poznata i provjerljiva, pa se pogreška može pripisati konkretnom
stupnju. \textbf{Mjerni etalon.} Castillo i suradnici koriste OptiTrack kao
komercijalnu referencu pri vrednovanju jeftinijih sustava praćenja, potvrđujući
ga kao točan izvor mjerenja \cite[str.~7]{castillo2025}.

\textbf{Filtriranje šumnog pokreta.} Martini i suradnici predlažu robustan
filtar za rad u stvarnom vremenu koji smanjuje podrhtavanje robota i uspoređuju
ga s linearnim Kalmanovim filtrom prvog i drugog reda
\cite[str.~5]{martini2024}, dok Zhu i suradnici naglašavaju da klasično
niskopropusno filtriranje unosi fazno kašnjenje nepovoljno za rad u stvarnom
vremenu te predlažu prilagodljivu atenuaciju tremora \cite[str.~2]{zhu2022}. To
opravdava izbor kauzalnog filtra One-Euro \cite{casiez2012} umjesto nekauzalnog
izglađivanja. Naši rezultati taj nalaz i zaoštravaju budući da se fazno kašnjenje filtra
pokazalo najvećim pojedinačnim doprinosom ukupnom odzivu sustava.

\textbf{Teleoperacija i autonomija.} Stroppa i suradnici razmatraju paradigme
dijeljenog upravljanja kao spektar između ručnog i autonomnog
\cite[str.~5]{stroppa2023}, a Glawe i suradnici sustavno pregledavaju ljudsku
autonomiju i osjećaj djelovanja (\emph{sense of agency}) u HRI-u
\cite[str.~1]{glawe2026}. Oba čine okosnicu rasprave u
poglavlju~\ref{sec:rasprava}.

\FloatBarrier
